% ============================================================================
% chapters/ch6_discussion.tex
% ============================================================================
\chapter{Discussion and Conclusion}
\label{ch:discussion}

% ============================================================================
\section{Answering the Research Questions}
\label{sec:rq_answers}
% ============================================================================

\paragraph{RQ1: Detecting degradation without labels.}
The three-layer monitoring framework detects a range of degradation scenarios
without any labelled data.
Confidence monitoring (Layer~1) detects a broad drop in prediction quality
including overconfident wrong predictions resulting from a unit mismatch.
Temporal consistency (Layer~2) detects erratic prediction sequences that arise
when sensor data is noisy or misaligned.
Distributional drift (Layer~3) detects shift in the input feature distribution
through PSI and z-score analysis.

The most sensitive proxy metric for actual accuracy was the per-channel
z-score (Layer~3): the accelerometer $a_z$ channel z-score of 312 during the
unit mismatch experiment provided unambiguous evidence of domain shift before
any accuracy measurement was possible.

\paragraph{RQ2: Proxy metrics and trigger thresholds.}
Multi-metric voting with two-of-three confirmation suppresses false positives
from transient noise while remaining sensitive to genuine drift.
The calibrated PSI thresholds (0.75 / 1.50) for multi-channel aggregated PSI
are a key empirical contribution: standard textbook thresholds (0.10 / 0.25)
were incorrectly sensitive, firing on stable data.

\paragraph{RQ3: AdaBN+TENT vs pseudo-labelling.}
Both strategies produced non-trivial adaptations.
AdaBN+TENT is faster (parameter-free AdaBN + 10 TENT gradient steps, no
labelled data from the source domain required) but resulted in a confidence
drop of 7.6 percentage points in the 2026-02-23 production run, failing the
canary gate.
Pseudo-labelling (A5) achieved 96.9\,\% in-loop validation accuracy but
requires a source-domain holdout for rollback evaluation, and is sensitive to
the self-consistency filter rate (only 14.5\,\% of source windows were
retained after filtering).

\paragraph{RQ4: Engineering safeguards.}
The experiments demonstrated that safeguards are essential.
Without the BN snapshot mechanism, TENT silently invalidated AdaBN's output.
Without the canary gate, the degraded AdaBN+TENT model (A4, 82.3\,\% mean
confidence) would have replaced the production model.
Without the governance flag \texttt{promote\_to\_shared=False}, a test run
would have overwritten the production baseline.

% ============================================================================
\section{Contributions in Context}
\label{sec:contributions_context}
% ============================================================================

The 14-stage composable pipeline architecture addresses the ``monolithic script''
problem identified by \citet{sculley2015hidden}, providing a modular system
where each stage is independently testable and composable.

The three-layer monitoring framework extends the single-metric monitoring
common in production ML literature to a composite health signal, reducing
false alarm rates through multi-metric voting.

The BN snapshot mechanism for the AdaBN+TENT chain is, to the author's
knowledge, not described in other implementations.
It resolves a fundamental interaction between two widely-used test-time
adaptation methods that would otherwise produce worse results than either
method alone.

% ============================================================================
\section{Limitations}
\label{sec:limitations}
% ============================================================================

\begin{enumerate}
  \item \textbf{Single-user evaluation.} The pipeline is validated on data
    from a small number of users.
    Generalisation to larger, more diverse populations is an open question.

  \item \textbf{No online ground truth.} The evaluation of monitoring quality
    relies on replaying known-bad scenarios (unit mismatch).
    In practice, the relationship between proxy metrics and true accuracy
    may vary across users and scenarios.

  \item \textbf{Pseudo-label quality ceiling.} When the model is already
    confused (low confidence, high entropy), pseudo-labels are unreliable.
    The curriculum strategy mitigates but does not eliminate this risk.

  \item \textbf{Label set fixed.} The pipeline handles only the 11
    pre-defined activity classes.
    Open-set recognition (detecting activities outside the training taxonomy)
    is not addressed.

  \item \textbf{Infrastructure dependencies.} The canary evaluation and
    rollback mechanisms require the MLflow tracking server to be running.
    If the server is unavailable, the pipeline falls back to a local
    registry file, losing audit trail completeness.
\end{enumerate}

% ============================================================================
\section{Future Work}
\label{sec:future_work}
% ============================================================================

\begin{description}[style=nextline]

  \item[Federated adaptation.]
    Extend the pseudo-labelling and AdaBN+TENT strategies to a federated
    setting where each user's device performs local adaptation and only
    updated model parameters are aggregated, preserving user data privacy.

  \item[Online learning.]
    Replace the batch-triggered retraining cycle with an incremental learning
    loop that updates the model continuously as new windows arrive, using
    reservoir sampling to maintain a representative memory buffer.

  \item[Automated PSI threshold calibration.]
    Implement a self-calibrating threshold mechanism that computes the
    99th-percentile PSI over a known-stable reference period and sets the
    warn threshold automatically, removing the need for manual calibration.

  \item[Open-set recognition.]
    Integrate an energy-based or prototype-based out-of-distribution detector
    that identifies activity sequences not belonging to any of the
    \nclasses{} training classes, enabling detection of new behaviours.

  \item[Prometheus / Grafana integration.]
    The Prometheus metrics exporter (\texttt{src/prometheus\_metrics.py})
    is implemented but not connected to a running Grafana instance in the
    current deployment.
    Completing this integration would enable real-time dashboard monitoring.

\end{description}

% ============================================================================
\section{Conclusion}
\label{sec:conclusion}
% ============================================================================

This thesis presented an end-to-end MLOps pipeline for continuous Human Activity
Recognition using wrist-worn inertial sensor data in an anxiety monitoring
context.
The pipeline addresses the central practical challenge of maintaining model
reliability in a setting where production data are unlabelled and sensor
characteristics may shift over time.

The core contributions --- a three-layer monitoring framework, a voting-based
trigger policy, the AdaBN+TENT BN snapshot mechanism, curriculum pseudo-labelling
with EWC regularisation, and a canary deployment gate --- together form a
coherent, safety-conscious system that can detect, diagnose, and recover from
domain shift without manual intervention.

The accelerometer unit mismatch case study demonstrates in concrete terms what
undetected domain shift looks like: 14.5\,\% cross-user accuracy with
paradoxically high confidence, fully reversible once the unit conversion
pipeline was introduced.
This case illustrates both the risk of unmonitored ML deployments and the
diagnostic value of the framework.

The pipeline is fully containerised, CI/CD-enabled, and reproducible.
The combination of Git, DVC, MLflow, Docker, and GitHub Actions provides the
provenance infrastructure needed to ensure that every model decision can be
traced to its data, code, and configuration state.

The highest-priority area for future improvement is the transition from
batch-triggered retraining to online continuous adaptation, which would reduce
response latency to domain shift from hours to seconds.
The composable stage architecture of the pipeline is designed to support this
transition incrementally.
